\documentclass[11pt]{article}
\usepackage[margin=1in]{geometry}
\usepackage{amsmath,amssymb,amsthm}
\usepackage{booktabs}
\usepackage[hidelinks]{hyperref}

\newtheorem{theorem}{Theorem}[section]
\newtheorem{proposition}[theorem]{Proposition}
\newtheorem{lemma}[theorem]{Lemma}
\newtheorem{corollary}[theorem]{Corollary}
\theoremstyle{remark}
\newtheorem{remark}[theorem]{Remark}

\newcommand{\Z}{\mathbb{Z}}
\newcommand{\TD}{\mathrm{TD}}
\newcommand{\RTD}{\mathrm{RTD}}

\title{The $4$-RGDD of type $10^{4}$ and three mutually orthogonal\\
Latin squares of order $10$:
an equivalence, what is known,\\ and a computer search}
\author{}
\date{}

\begin{document}
\maketitle

\begin{abstract}
In the existence theorems of Wei and Ge for $4$-RGDDs
\cite[Theorem~4.21]{WG14}, \cite[Theorem~5.2]{WG17} the type $10^{4}$ is one of the
possible exceptions.  We were asked to construct a $4$-RGDD of type $10^{4}$
directly, i.e.\ three mutually orthogonal Latin squares of order $10$, or else to
prove that none exists.
\textbf{We achieved neither.}
We prove that a $4$-RGDD of type $10^{4}$ exists if and only if three mutually
orthogonal Latin squares of order $10$ exist (Theorem~\ref{thm:equiv}); the
existence of three MOLS$(10)$ is a famous open problem, stated as unresolved in work
published in 2026 \cite{BKS26}.  A construction would settle it one way and a
non-existence proof the other, so neither can be expected from routine methods.
What this note does contain:
(i) complete proofs of the equivalence and of two obstructions to the usual direct
constructions (no $4$-GDD of type $10^{4}$ admits a group of order $10$ fixing three of
its groups and acting regularly on them; in particular none admits an abelian
point-regular group, so none is cyclic);
(ii) a verified incomplete $4$-RGDD of type $10^{4}$ with a hole of type $2^{4}$,
obtained from Brouwer's four MOLS of order $10$ with a hole of order $2$, and a proof
that filling holes cannot help: a $4$-IRGDD of type $(10,h)^{4}$ has $h\le 2$, for $h=2$
the hole can never be filled, and $h=1$ is equivalent to the original problem (in
particular a hole of size $4$, which three MOLS of order $4$ could fill, cannot occur: there
are not even two incomplete MOLS of order $10$ with a common $4\times 4$ hole);
(iii) an independent verification of the near triple of Egan and Wanless, which in
design language is a $4$-GDD of type $10^{4}$ with seven pairwise disjoint parallel
classes (a resolution needs ten);
(iv) a computer search in which \textbf{$198{,}464$} Latin squares of order $10$ were each
tested \emph{exhaustively} for membership in a triple of MOLS; none is.
\end{abstract}

\section{Introduction}

The existence theorem for $4$-RGDDs, in the form updated by Wei and Ge
\cite[Theorem~4.21]{WG14}, \cite[Theorem~5.2]{WG17}, states that the
necessary conditions for a $4$-RGDD of type $h^{u}$,
\[
u\ge 4,\qquad hu\equiv 0 \pmod 4,\qquad h(u-1)\equiv 0\pmod 3,
\]
are sufficient except for $(h,u)\in\{(2,4),(2,10),(3,4),(6,4)\}$ and a finite list of
possible exceptions; for $h=10$ the possible exceptions are $u\in\{4,94\}$
\cite[Theorem~5.2]{WG17}.  A $4$-RGDD of type $10^{4}$ has
\[
v=40 \text{ points},\qquad b=100 \text{ blocks},\qquad r=10 \text{ parallel classes of } 10 \text{ blocks each},
\]
and it satisfies the necessary conditions ($40\equiv 0 \pmod 4$, $30\equiv 0\pmod 3$).

\medskip\noindent
\textbf{The task and its outcome.}
The request was to construct such a design directly --- equivalently, three pairwise
orthogonal Latin squares of order $10$ --- and, if that could not be done, to give a
proof.  \textbf{No such design is constructed here, and no proof of non-existence is
given.}  By Theorem~\ref{thm:equiv} either one would decide whether $N(10)\ge 3$,
one of the best-known open problems about Latin squares
(``It is unknown if $N(10)\ge 3$, and determining the value of $N(10)$ is one of the
most prominent unsolved problems concerning MOLS'' \cite[Section~1]{BKS26}).
We do not present heuristic arguments as proofs.  What can be proved, and what we
found, is organised as follows:
Section~\ref{sec:equiv} proves the equivalence;
Section~\ref{sec:known} summarises what is known about three MOLS$(10)$ and explains why
we cannot supply a proof either way;
Section~\ref{sec:obstr} proves two obstructions relevant to direct constructions;
Section~\ref{sec:near} gives explicit, verified near misses;
Section~\ref{sec:search} describes our computer search.
The script \texttt{verify\_10\_4.py} checks every object printed here from the
definitions, and can check any candidate triple of squares.

\section{Definitions and the equivalence}\label{sec:equiv}

Throughout $[g]=\{0,1,\dots,g-1\}$.
A \emph{$k$-GDD of type $g^{u}$} is a triple $(X,\mathcal G,\mathcal B)$ where $X$ is a set
of $gu$ points, $\mathcal G$ is a partition of $X$ into $u$ \emph{groups} of size $g$,
and $\mathcal B$ is a set of $k$-subsets of $X$ (\emph{blocks}) such that every pair of
points from distinct groups lies in exactly one block and no block contains two points
of the same group.  A \emph{parallel class} is a set of blocks partitioning $X$; the
GDD is \emph{resolvable} (a $k$-RGDD) if $\mathcal B$ can be partitioned into parallel
classes.  A $\TD(k,g)$ is a $k$-GDD of type $g^{k}$, and an $\RTD(k,g)$ is a resolvable
one.

A \emph{Latin square} of order $g$ is a $g\times g$ array with entries in $[g]$ in which
every symbol occurs once in each row and once in each column.  Latin squares $A$ and $B$
are \emph{orthogonal} ($A\perp B$) if the $g^{2}$ ordered pairs $(A(i,j),B(i,j))$ are
distinct; MOLS are pairwise orthogonal Latin squares, and $N(g)$ is the largest number
of MOLS of order $g$.  A \emph{transversal} of $A$ is a set of $g$ cells containing
one cell of each row and each column and each symbol once; a \emph{common transversal}
of $A$ and $B$ is a transversal of both.  If $A\perp B$ then each symbol class
$\{(i,j):B(i,j)=s\}$ is a transversal of $A$; conversely every partition of the cells
into $g$ transversals of $A$ defines (after labelling the transversals by $[g]$) a
Latin square orthogonal to $A$, an \emph{orthogonal mate} of $A$.

\begin{lemma}\label{lem:td}
In a $4$-GDD of type $g^{4}$ (that is, a $\TD(4,g)$) every block meets every group in
exactly one point and there are exactly $g^{2}$ blocks.  Every parallel class has $g$
blocks, and a resolution has $g$ parallel classes.
\end{lemma}

\begin{proof}
A block has four points, no two in one group, and there are four groups.  Each block
contains $\binom42=6$ pairs from distinct groups; there are $6g^{2}$ such pairs and each
lies in exactly one block.  A parallel class covers the $4g$ points by blocks of size $4$.
\end{proof}

\begin{theorem}\label{thm:equiv}
Let $g\ge 2$.  A $4$-RGDD of type $g^{4}$ exists if and only if there are three MOLS of
order $g$, if and only if a $\TD(5,g)$ exists.  Explicitly: let $A,B,C$ be MOLS of order
$g$, let
\[
X=\{r_i,\;c_i,\;x_i,\;y_i : i\in[g]\}
\]
with the four groups $\{r_i\}$, $\{c_i\}$, $\{x_i\}$, $\{y_i\}$, and for every cell
$(i,j)$ let
\[
\beta_{ij}=\{\,r_i,\;c_j,\;x_{A(i,j)},\;y_{B(i,j)}\,\}.
\]
Then the blocks $\beta_{ij}$ form a $4$-GDD of type $g^{4}$, and the sets
$\mathcal P_s=\{\beta_{ij}: C(i,j)=s\}$, $s\in[g]$, form a resolution of it.  Every
$4$-RGDD of type $g^{4}$ arises in this way.
\end{theorem}

\begin{proof}
\emph{From squares to the design.}  The pair $\{r_i,c_j\}$ lies only in $\beta_{ij}$;
the pair $\{r_i,x_s\}$ lies in $\beta_{ij}$ exactly when $A(i,j)=s$, which happens for
exactly one $j$ because row $i$ of $A$ contains $s$ once; the same argument with
columns and with $B$ handles $\{c_j,x_s\}$, $\{r_i,y_t\}$, $\{c_j,y_t\}$; and
$\{x_s,y_t\}$ lies in $\beta_{ij}$ exactly when $(A(i,j),B(i,j))=(s,t)$, which happens
once because $A\perp B$.  No block contains two points of a group.  The cells with
$C(i,j)=s$ form a common transversal of $A$ and $B$: one cell in each row and column
because $C$ is Latin, and each symbol of $A$ (of $B$) once because $C\perp A$
($C\perp B$).  Hence the blocks of $\mathcal P_s$ contain each $r_i$, $c_j$, $x_s$, $y_t$
exactly once, i.e.\ $\mathcal P_s$ is a parallel class, and the $\mathcal P_s$ partition
the blocks.

\emph{From the design to squares.}  Let $(X,\mathcal G,\mathcal B)$ be a $4$-RGDD of type
$g^{4}$; name the groups and their points $\{r_i\},\{c_j\},\{x_s\},\{y_t\}$ and the
parallel classes $\mathcal P_0,\dots,\mathcal P_{g-1}$ (Lemma~\ref{lem:td}).  By
Lemma~\ref{lem:td} the unique block through $r_i$ and $c_j$ has the form
$\{r_i,c_j,x_s,y_t\}$; put $A(i,j)=s$, $B(i,j)=t$ and let $C(i,j)$ be the index of the
parallel class containing that block.  The pair conditions for $\{r_i,x_s\}$ and
$\{c_j,x_s\}$ say that $A$ is Latin, likewise for $B$, and the condition for
$\{x_s,y_t\}$ says $A\perp B$.  A parallel class contains each $r_i$ and each $c_j$
once, so $C$ is Latin, and it contains each $x_s$ and each $y_t$ once, so the pairs
$(A,C)$ and $(B,C)$ take every value once: $C\perp A$ and $C\perp B$.

\emph{$\TD(5,g)$.}  Given a resolution $\mathcal P_0,\dots,\mathcal P_{g-1}$, adjoin a fifth
group $\{\infty_0,\dots,\infty_{g-1}\}$ and add $\infty_s$ to every block of $\mathcal P_s$;
since $\mathcal P_s$ partitions $X$, every pair $\{z,\infty_s\}$ lies in exactly one block,
so this is a $\TD(5,g)$.  Conversely, in a $\TD(5,g)$ the blocks through a point $\infty_s$
of the fifth group, with $\infty_s$ removed, form a parallel class of the $\TD(4,g)$ obtained
by deleting that group, and these classes partition its blocks.
\end{proof}

\begin{corollary}\label{cor:types}
For $g\ge 2$ a $4$-RGDD of type $g^{4}$ exists if and only if $N(g)\ge 3$.  As $N(2)=1$,
$N(3)=2$, $N(6)=1$ (Tarry \cite{Tar1900}) and $N(g)\ge 3$ for every other $g\ne 10$
(\cite{ACD07}; consistently, \cite[Theorem~5.2]{WG17} lists $g=10$ as the only possible
exception of type $g^{4}$), such a design exists
for all $g\notin\{2,3,6,10\}$ and does not exist for
$g\in\{2,3,6\}$ --- exactly the pattern of exceptions $(2,4),(3,4),(6,4)$ and possible
exception $(10,4)$ in \cite{WG14,WG17}.  For $g=10$ the question is equivalent to $N(10)\ge 3$.
\end{corollary}

\begin{corollary}\label{cor:gdd}
Every $4$-GDD of type $10^{4}$ is the design $\{\beta_{ij}\}$ of a pair of orthogonal
Latin squares $(A,B)$ of order $10$ (such pairs exist \cite{Par59,BSP60}), and its
parallel classes are exactly the sets $\{\beta_{ij}:(i,j)\in T\}$ with $T$ a common
transversal of $A$ and $B$.  Hence a $4$-GDD of type $10^{4}$ is resolvable if and only
if the $100$ cells can be partitioned into $10$ common transversals of $A$ and $B$.
\end{corollary}

\begin{proof}
The first statement is the second part of the proof of Theorem~\ref{thm:equiv}, which
does not use the resolution to define $A$ and $B$.  The blocks over a set $T$ of cells
contain every point exactly once if and only if $T$ is a common transversal.
\end{proof}

So $4$-GDDs of type $10^{4}$ exist in abundance; the entire difficulty is resolvability.

\section{What is known, and why we cannot give a proof}\label{sec:known}

\paragraph{Bounds.}  $N(10)\ge 2$: Parker \cite{Par59} found two orthogonal Latin squares
of order $10$ (see also \cite{BS59,BSP60}).  $N(10)\le 6$: there is no projective
plane of order $10$ \cite{LTS89} (see also the independent searches
\cite{Roy11,BCSKG21}), so $N(10)<9$, and together with a completion result for nets due
to Bruck \cite{Bru63} this gives $N(10)\le 6$, the best known upper bound
\cite[Section~2.4]{BKS26}, \cite{DH14}.

\paragraph{Partial non-existence results.}  Each of the following excludes triples
of MOLS$(10)$ with a special structure (statements as summarised in \cite{BKS26}).
\begin{itemize}
\item Mann \cite{Man44}: a Latin square of order $10$ with a $5\times 5$ Latin subsquare
  has no orthogonal mate (Proposition~\ref{prop:mann} below).
\item Parker \cite{Par62}: two orthogonal Latin squares of order $10$ with orthogonal
  $3\times 3$ Latin subsquares cannot be part of an orthogonal triple.
\item McKay, Meynert and Myrvold \cite{MMM07}: every Latin square in a triple of
  MOLS$(10)$ has trivial symmetry group.
\item Myrvold \cite{Myr99}, Bright, Keita and Stevens \cite{BKS26}: if one square of a
  triple has a $4\times 4$ Latin subsquare, the other two fall into one of $28$
  ``mate pattern type pairs''; Myrvold excluded $20$ of them, and \cite{BKS26}
  re-verified this with a SAT solver and found orthogonal pairs in all eight remaining
  cases, so those cases cannot be excluded by looking at pairs alone.
\item Gill and Wanless \cite{GW23}: the rank of the code generated by any pair of squares
  of a triple of MOLS$(10)$ is exactly $37$; see also Delisle \cite{Del10} and Dukes and
  Howard \cite{DH14}.
\item Zaikin, Vatutin and Bright \cite{ZVB25}: no extended self-orthogonal diagonal Latin
  square of order $10$ lies in an orthogonal triple.
\end{itemize}

\paragraph{Near misses.}  Brouwer \cite{Bro84}: four MOLS of order $10$ with a common
hole of order $2$, i.e.\ a $\TD(6,10)-\TD(6,2)$ (Section~\ref{sec:brouwer}).
Egan and Wanless \cite{EW16}: squares $L_1,L_2,L_3$ of order $10$ with
$L_1\perp L_2$, $L_1\perp L_3$ and $91$ of the $100$ ordered pairs distinct when
$L_2$ and $L_3$ are superimposed; hence $L_1$ and $L_2$ have seven pairwise disjoint common transversals, improving
the four of Brown, Hedayat and Parker \cite{BHP93} (Section~\ref{sec:ew}).  Zaikin and Kochemazov \cite{ZK15}: three
diagonal Latin squares of order $10$ satisfying the orthogonality condition on a common
set of $73$ cells.

\paragraph{Why we cannot give a proof.}
\begin{enumerate}
\item By Theorem~\ref{thm:equiv}, a proof that no $4$-RGDD of type $10^{4}$ exists is a
  proof that $N(10)=2$, and a construction proves $N(10)\ge 3$.  Either would settle the
  open problem above.  We have neither.
\item All the standard necessary conditions are satisfied: the divisibility conditions
  hold, the underlying $4$-GDD of type $10^{4}$ exists (Corollary~\ref{cor:gdd}), and so
  do incomplete versions with a small hole (Section~\ref{sec:brouwer}).  Non-existence
  therefore cannot follow from these conditions.
\item The proof that $N(10)<9$ \cite{LTS89} relies on the binary code of a projective
  plane and still needed a very large computation.  The methods behind the partial
  results above (symmetry, subsquares, codes of pairs) each exclude only special
  configurations.
\item A naive exhaustive search would test one square from each of the
  $34{,}817{,}397{,}894{,}749{,}939$ main classes of Latin squares of order $10$
  \cite{MMM07}.  (Membership in a triple is a main-class invariant: a square is a
  component of a triple if and only if its $\TD(3,10)$ extends to a $\TD(5,10)$, and this
  does not depend on which of its three groups are called rows, columns and symbols.)
  At the speed of our own complete test (about $27$ squares per second per core,
  Section~\ref{sec:search}) this would take about $4.1\times 10^{7}$ core-years.
\end{enumerate}

\section{Two obstructions with complete proofs}\label{sec:obstr}

\begin{proposition}[Mann \cite{Man44}]\label{prop:mann}
A Latin square of order $4t+2$ containing a Latin subsquare of order $2t+1$ has no
transversal.  In particular no Latin square of order $10$ with a $5\times 5$ Latin
subsquare has an orthogonal mate, so none can be one of the squares $A,B,C$ of
Theorem~\ref{thm:equiv}.
\end{proposition}

\begin{proof}
Let the subsquare occupy rows $R$ and columns $K$ with symbol set $S$,
$|R|=|K|=|S|=m=2t+1$, and let bars denote complements ($|\bar R|=|\bar K|=|\bar S|=m$).
Rows in $R$ carry the symbols of $S$ in the columns $K$, hence the symbols of
$\bar S$ in the columns $\bar K$; likewise columns in $K$ carry $\bar S$ in the rows
$\bar R$; consequently the cells $\bar R\times\bar K$ carry only symbols of $S$.  Let $T$
be a transversal and let $a,b,d$ be the numbers of its cells in $R\times K$,
$R\times\bar K$, $\bar R\times\bar K$.  Counting the cells of $T$ in the rows $R$, in the
columns $\bar K$, and those with a symbol in $S$ gives
\[
a+b=m,\qquad b+d=m,\qquad a+d=m .
\]
The first and third give $b=d$, and then $2b=m$, which is impossible since $m$ is odd.
\end{proof}

A \emph{complete mapping} of a group $H$ is a bijection $\theta\colon H\to H$ such that
$x\mapsto x\theta(x)$ is also a bijection; equivalently the cells $(x,\theta(x))$ form a
transversal of the Cayley table $(x,y)\mapsto xy$.

\begin{proposition}\label{prop:cm}
No group of order $10$ has a complete mapping.
\end{proposition}

\begin{proof}
A group $H$ of order $10$ has a normal subgroup $N$ of order $5$ (its Sylow $5$-subgroup);
let $\varepsilon\colon H\to H/N\cong\Z_2$, written additively.  If $\theta$ were a complete
mapping, then summing $\varepsilon(x\theta(x))=\varepsilon(x)+\varepsilon(\theta(x))$ over
$x\in H$ would give $\sum_{y\in H}\varepsilon(y)=2\sum_{y\in H}\varepsilon(y)$ in $\Z_2$,
because $x\mapsto x\theta(x)$ and $\theta$ are bijections.  But
$\sum_{y\in H}\varepsilon(y)=|H\setminus N|=5\equiv 1$, a contradiction.  (This is a special
case of the Hall--Paige theorem \cite{HP55}.)
\end{proof}

\begin{lemma}\label{lem:dm}
Let $(X,\{G_1,G_2,G_3\},\mathcal B)$ be a $\TD(3,m)$ and let $H$ be a group of order $m$
acting on it by automorphisms that map each $G_i$ onto itself and act regularly on each
$G_i$.  Then $H$ has a complete mapping.
\end{lemma}

\begin{proof}
Fix $p_i\in G_i$ and write $(h,i)$ for the point $h(p_i)$; so $g\in H$ maps $(h,i)$ to
$(gh,i)$.  For $a\in H$ let $\{(1,1),(a,2),(\theta(a),3)\}$ be the block through $(1,1)$ and
$(a,2)$.  The map $\theta$ is injective, because the pair $\{(1,1),(\theta(a),3)\}$ lies in
only one block, hence bijective.  Every block contains exactly one point $(h,1)$, and
$h^{-1}$ maps it to a block through $(1,1)$; so the blocks are exactly
$\{(h,1),(ha,2),(h\theta(a),3)\}$ for $h,a\in H$.  The pair $\{(y,2),(z,3)\}$ lies in the
block with parameters $(h,a)$ if and only if $h=ya^{-1}$ and $a^{-1}\theta(a)=y^{-1}z$.
Since every such pair lies in exactly one block, $a\mapsto a^{-1}\theta(a)$ is a
bijection.  With $\sigma(x)=\theta(x^{-1})$ we get $x\sigma(x)=a^{-1}\theta(a)$ for
$a=x^{-1}$, so $\sigma$ is a complete mapping.
\end{proof}

\begin{corollary}\label{cor:aut}
\begin{enumerate}
\item[(a)] No $4$-GDD of type $10^{4}$, resolvable or not, admits a group $H$ of order $10$
  of automorphisms that maps three of its groups onto themselves and acts regularly on
  each of them.  In particular no $\TD(k,10)$ with $k\ge 3$, and so no $4$-RGDD of type
  $10^{4}$, comes from a $(10,k;1)$-difference matrix over a group of order $10$
  (compare \cite[Lemma~2.8]{WG14}, which gives such transversal designs with $k=2$ only).
\item[(b)] No $4$-GDD of type $10^{4}$ admits an abelian group of automorphisms acting
  regularly on its $40$ points.  In particular none is cyclic, and none is obtained by
  developing base blocks under a regular action of one of the abelian groups of order
  $40$, namely $\Z_{40}$, $\Z_{20}\times\Z_2\cong\Z_{10}\times\Z_4$ and
  $\Z_{10}\times\Z_2^{2}$.
\item[(c)] If $G$ is any group of automorphisms of a $4$-GDD of type $10^{4}$ acting
  regularly on the $40$ points, then the kernel $K$ of the action of $G$ on the four
  groups has order $5$ and $G/K$ is dihedral of order $8$, acting on the four groups as
  the symmetry group of a square.
\end{enumerate}
\end{corollary}

\begin{proof}
An automorphism maps groups onto groups, since two points lie in a common group if and
only if no block contains both.
(a) Deleting a fourth group gives a $\TD(3,10)$ to which Lemma~\ref{lem:dm} applies,
contradicting Proposition~\ref{prop:cm}.  A $(10,k;1)$-difference matrix over $H$
yields a $\TD(k,10)$ on which $H$ acts by translation, fixing every group and regularly
on each.
(b), (c) Let $G$ act regularly; then it acts transitively on the four groups.  Let
$\Gamma$ be the group containing a point $p_0$ and $H$ its stabiliser, so $|H|=10$, $H$
acts regularly on $\Gamma$, and the kernel $K$ of the action on the groups is the
largest normal subgroup of $G$ contained in $H$.  Now $G/K$ embeds in $S_4$ and is
transitive on four objects, so $|G/K|$ is a multiple of $4$ dividing $\gcd(40,24)=8$.  If
$|G/K|=4$ then $K=H$ fixes every group and acts on each freely, hence regularly, which
(a) forbids.  So $|K|=5$ and $|G/K|=8$; the only subgroups of order $8$ of $S_4$ are its
Sylow $2$-subgroups, dihedral of order $8$.  If $G$ is abelian then $H$ is normal, so
$K=H$, which is impossible; this proves (b).
\end{proof}

\begin{remark}\label{rem:mmm}
Combined with the computer-assisted theorem of \cite{MMM07} quoted in
Section~\ref{sec:known} (every square of a triple of MOLS$(10)$ has trivial symmetry
group), much more follows.  Let $\varphi\ne 1$ be an automorphism of the $4$-GDD
underlying a $4$-RGDD of type $10^{4}$ (the resolution need not be preserved) that maps
each of the four groups onto itself.  Then $\varphi$ permutes the rows,
the columns and the symbols of the square $A$ of Theorem~\ref{thm:equiv} and maps the
block of cell $(i,j)$ to the block of the image cell, so it induces an autotopism of $A$;
this autotopism is non-trivial, since otherwise $\varphi$ would fix every block through
every pair $r_i,c_j$ and hence every point.  As $A$ is a component of a triple, this
contradicts \cite{MMM07}.  So \emph{every non-identity automorphism of a $4$-RGDD of type
$10^{4}$ moves at least one of its groups}; in particular the subgroup $K$ in (c) cannot
exist, and \emph{no $4$-RGDD of type $10^{4}$ admits a point-regular automorphism group}.
These two statements rest on the computation of \cite{MMM07}, which we have not repeated.
\end{remark}

\section{Explicit near misses (verified)}\label{sec:near}

\subsection{\texorpdfstring{An incomplete $4$-RGDD of type $10^{4}$ with a hole of type $2^{4}$}{An incomplete 4-RGDD of type 10\textasciicircum 4 with a hole of type 2\textasciicircum 4}}\label{sec:brouwer}

Brouwer \cite{Bro84} gave four MOLS of order $10$ with a common hole of order $2$
(Figure~\ref{fig:brouwer}, transcribed from \cite[Fig.~1]{Bro84}; entry $ijkl$ in cell
$(r,c)$ means that the four squares have symbols $i,j,k,l$ there; rows and columns are
numbered $0,\dots,9$, the hole is $\{8,9\}\times\{8,9\}$, and the hole symbols are $8,9$).

\begin{figure}[ht]
\centering
\begin{tabular}{*{10}{c}}
5039 & 8225 & 4568 & 9753 & 6812 & 0390 & 2487 & 7976 & 3141 & 1604\\
9460 & 6109 & 8336 & 5678 & 4947 & 7823 & 1091 & 3584 & 0252 & 2715\\
6748 & 9571 & 7219 & 8007 & 0685 & 5954 & 4830 & 2192 & 1363 & 3426\\
8114 & 7458 & 9642 & 4329 & 3293 & 1786 & 6965 & 5801 & 2070 & 0537\\
2856 & 4794 & 6083 & 3932 & 1479 & 8661 & 0128 & 9317 & 7505 & 5240\\
0903 & 3867 & 5495 & 7180 & 9024 & 2549 & 8772 & 1238 & 4616 & 6351\\
4281 & 1910 & 0874 & 6596 & 2308 & 9135 & 3659 & 8443 & 5727 & 7062\\
7697 & 5382 & 2921 & 1845 & 8550 & 3018 & 9206 & 0769 & 6434 & 4173\\
1522 & 2633 & 3700 & 0411 & 5166 & 6277 & 7344 & 4055 & --   & --  \\
3375 & 0046 & 1157 & 2264 & 7731 & 4402 & 5513 & 6620 & --   & --
\end{tabular}
\caption{Four MOLS of order $10$ with a hole of order $2$ (Brouwer \cite{Bro84}).}
\label{fig:brouwer}
\end{figure}

\begin{proposition}\label{prop:irgdd}
Let $A,B,C$ be the first three of the squares in Figure~\ref{fig:brouwer}, defined on the
$96$ cells outside the hole.  With $X$ and $\beta_{ij}$ as in Theorem~\ref{thm:equiv} and
\[
Y=\{r_8,r_9,c_8,c_9,x_8,x_9,y_8,y_9\},
\]
the $96$ blocks $\beta_{ij}$ have the following properties.
\begin{enumerate}
\item[(i)] Every pair of points from distinct groups that are not both in $Y$ lies in
  exactly one block, no pair inside $Y$ lies in a block, and every block meets $Y$ in at
  most one point.  (This is a $4$-IGDD of type $(10,2)^{4}$ in the notation of
  \cite{WG14}.)
\item[(ii)] For $s\in\{0,\dots,7\}$ the blocks of the cells with $C(i,j)=s$ form a parallel
  class ($10$ blocks); for $s\in\{8,9\}$ they form a holey parallel class ($8$ blocks
  partitioning $X\setminus Y$).  So the design is a $4$-IRGDD of type $(10,2)^{4}$.
\item[(iii)] The hole cannot be filled: there is no set of blocks of size $4$ which,
  added to these $96$ blocks, yields a $4$-GDD of type $10^{4}$, let alone a $4$-RGDD.
\end{enumerate}
\end{proposition}

\begin{proof}
(i) and (ii) say that $A,B,C$ are incomplete MOLS with a common hole (every row and
column outside the hole contains every symbol once, the rows and columns of the hole
contain each of $0,\dots,7$ once outside the hole, and each pair of squares shows every
ordered pair of symbols not both in $\{8,9\}$ exactly once); the argument of
Theorem~\ref{thm:equiv} translates this into (i) and (ii).  All of this is checked by
\texttt{verify\_10\_4.py}, which in fact verifies that all four squares form a
$\TD(6,10)-\TD(6,2)$.
(iii) The pairs not yet covered are exactly the $24$ pairs of points of $Y$ in distinct
groups.  An added block may contain no pair that is already covered, so all its points
lie in $Y$; hence the added blocks would form a $4$-GDD of type $2^{4}$ on $Y$.  By the
second part of the proof of Theorem~\ref{thm:equiv} (with $g=2$) such a design gives two
orthogonal Latin squares of order $2$, which do not exist (see also \cite{BSH77}).
\end{proof}

Thus a resolvable design on $40$ points that agrees with a $4$-RGDD of type $10^{4}$ on
all but the $4$ blocks inside an $8$-point hole exists, and the standard ``filling in
holes'' construction fails here only because $N(2)=1$.  No other incomplete design can
do better:

\begin{proposition}\label{prop:holes}
Let $1\le h\le 9$ and suppose that a $4$-IRGDD of type $(10,h)^{4}$ exists, with hole $Y$
($|Y|=4h$).  Then:
\begin{enumerate}
\item[(i)] $h\le 2$;
\item[(ii)] if $h=2$, no set of blocks can be added to obtain a $4$-GDD of type $10^{4}$;
\item[(iii)] if $h=1$, adding $Y$ as a block gives a $4$-RGDD of type $10^{4}$.
\end{enumerate}
So the ``filling in holes'' approach can only produce a $4$-RGDD of type $10^{4}$ from a
design that is already one with a block removed.
\end{proposition}

\begin{proof}
A point of $Y$ forms a pair with each of the $3(10-h)$ points outside $Y$ and outside its
own group, and every block through it contains three of these (a block meets $Y$ at most
once); so it lies in $10-h\ge 1$ blocks, all of them in (full) parallel classes, since
the holey ones avoid $Y$.  Likewise a point outside $Y$ lies in $10$ blocks, one in each
parallel class and one in each holey parallel class; hence there are $10-h$ parallel and
$h$ holey parallel classes.  A parallel class has $10$ blocks, each meeting $Y$ in at most
one point, and it covers the $4h$ points of $Y$; so $4h\le 10$, which is (i).
(ii) is the argument of Proposition~\ref{prop:irgdd}(iii), which used only the hole.
(iii) If $h=1$ the uncovered pairs are the six pairs inside the $4$-set $Y$, which meets
every group once, so $Y$ can be added as a block; added to the unique holey parallel
class, it completes a parallel class.
\end{proof}

\begin{remark}[Holes of size $h\ge 3$, in particular $h=4$]\label{rem:h4}
For $h\in\{4,5,7,8,9\}$ the design needed to fill the hole does exist: $N(h)\ge 3$, so by
Theorem~\ref{thm:equiv} there is a $4$-RGDD of type $h^{4}$; for $h=4$ it comes from three
MOLS of order $4$.  The obstruction is therefore the incomplete design, not the filling.
For $h=4$ the hole $Y$ has $16$ points, and every point of $Y$ lies in $10-h=6$ blocks, all
in full parallel classes; but a full parallel class has only $10$ blocks, each with at
most one point of $Y$, so it cannot cover the $16$ points of $Y$.

Even the non-resolvable version fails: there are not even two incomplete MOLS of order $10$
with a common $4\times 4$ hole (a $\TD(4,10)-\TD(4,4)$).  Let $L,M$ be such squares, with
hole rows $R$, hole columns $K$ and hole symbols $S$, $|R|=|K|=|S|=4$.  A hole symbol never
occurs in a hole row or a hole column, so in each of the six rows outside $R$ the four hole
symbols of $L$ lie in the six columns outside $K$: $L$ has a hole symbol in $24$ of the
$36$ cells of $\bar R\times\bar K$, and so has $M$.  Hence at least $24+24-36=12$ cells
carry a pair of hole symbols, which an incomplete pair must not contain.  The same count
gives the standard bound: $k$ incomplete MOLS of order $n$ with a common hole of size $h$
($0<h<n$) need $k\,h(n-h)\le (n-h)^{2}$, i.e.\ $n\ge (k+1)h$.  For $n=10$ this means
$h\le 3$ for a pair and $h\le 2$ for a triple.

The bound also disposes of the variants with several filled sub-designs, as in
\cite[Construction~4.3]{WG14} with $10=h_1+h_2+\cdots$: removing a filled sub-design of
type $h_i^{4}$, whose parallel classes extend to parallel classes of the whole design,
leaves a $4$-IRGDD of type $(10,h_i)^{4}$.  So every $h_i\le 2$, a part $h_i=2$ cannot be
filled, and only $h_i=1$ (single blocks) remains.

As an independent check, a SAT solver (CaDiCaL via PySAT, script \texttt{imols\_sat.py})
reports within a second that no two and no three incomplete MOLS of order $10$ with a
common $4\times 4$ hole exist, and no three with a common $3\times 3$ hole; on the other
hand it finds at once an incomplete Latin square of order $10$ with a $4\times 4$ hole,
two incomplete MOLS of order $12$ with a common $4\times 4$ hole, and two incomplete MOLS
of order $9$ with a common $3\times 3$ hole (all solutions checked by the script).
\end{remark}

\subsection{The near triple of Egan and Wanless}\label{sec:ew}

Egan and Wanless \cite[Section~8]{EW16} examined the Latin squares of order $10$ with an
autoparatopism group of order at least $3$ (catalogued in \cite{MMM07}) and found the
squares $A,B,C$ of Figure~\ref{fig:ew}, which are closer to a triple of MOLS than any
previously published example.

\begin{figure}[ht]
\centering
\setlength{\arraycolsep}{2.6pt}
$\begin{array}{*{10}{c}@{\hspace{2em}}*{10}{c}@{\hspace{2em}}*{10}{c}}
0 & 8 & 9 & 7 & 5 & 6 & 4 & 2 & 3 & 1 & 0 & 7 & 8 & 9 & 1 & 2 & 3 & 4 & 5 & 6 & 0 & 7 & 8 & 9 & 1 & 2 & 3 & 4 & 5 & 6\\
9 & 1 & 4 & 6 & 2 & 7 & 3 & 8 & 0 & 5 & 9 & 0 & 6 & 1 & 8 & 3 & 2 & 5 & 4 & 7 & 6 & 4 & 2 & 8 & 9 & 5 & 1 & 3 & 7 & 0\\
7 & 4 & 2 & 5 & 1 & 3 & 8 & 6 & 9 & 0 & 7 & 2 & 0 & 4 & 3 & 9 & 1 & 8 & 6 & 5 & 4 & 9 & 5 & 3 & 2 & 7 & 6 & 0 & 1 & 8\\
8 & 6 & 5 & 3 & 9 & 2 & 1 & 0 & 4 & 7 & 8 & 5 & 3 & 0 & 2 & 1 & 7 & 6 & 9 & 4 & 5 & 1 & 7 & 6 & 4 & 3 & 8 & 9 & 0 & 2\\
6 & 2 & 1 & 8 & 4 & 0 & 9 & 5 & 7 & 3 & 6 & 9 & 5 & 3 & 0 & 7 & 4 & 2 & 1 & 8 & 3 & 2 & 9 & 0 & 7 & 1 & 5 & 6 & 8 & 4\\
4 & 9 & 3 & 2 & 7 & 5 & 0 & 1 & 6 & 8 & 4 & 1 & 7 & 6 & 5 & 0 & 8 & 9 & 3 & 2 & 1 & 0 & 3 & 7 & 6 & 8 & 2 & 5 & 4 & 9\\
5 & 3 & 7 & 1 & 0 & 8 & 6 & 9 & 2 & 4 & 5 & 4 & 2 & 8 & 9 & 6 & 0 & 3 & 7 & 1 & 2 & 8 & 0 & 1 & 3 & 4 & 9 & 7 & 6 & 5\\
3 & 5 & 0 & 9 & 8 & 4 & 2 & 7 & 1 & 6 & 3 & 6 & 1 & 7 & 4 & 8 & 5 & 0 & 2 & 9 & 9 & 5 & 4 & 2 & 8 & 6 & 0 & 1 & 3 & 7\\
1 & 7 & 6 & 0 & 3 & 9 & 5 & 4 & 8 & 2 & 1 & 8 & 4 & 2 & 6 & 5 & 9 & 7 & 0 & 3 & 7 & 3 & 6 & 5 & 0 & 9 & 4 & 8 & 2 & 1\\
2 & 0 & 8 & 4 & 6 & 1 & 7 & 3 & 5 & 9 & 2 & 3 & 9 & 5 & 7 & 4 & 6 & 1 & 8 & 0 & 8 & 6 & 1 & 4 & 5 & 0 & 7 & 2 & 9 & 3\\
\end{array}$
\caption{The squares $A$ (left), $B$ (middle), $C$ (right) of Egan and Wanless
\cite[Section~8]{EW16}.}
\label{fig:ew}
\end{figure}

\begin{proposition}\label{prop:ew}
For the squares of Figure~\ref{fig:ew}:
\begin{enumerate}
\item[(i)] $A\perp B$, $A\perp C$, and superimposing $B$ and $C$ gives $91$ distinct
  ordered pairs; every repeated pair has $C$-symbol $7$, $8$ or $9$ \cite{EW16}.
\item[(ii)] Consequently the $4$-GDD of type $10^4$ of the pair $(A,B)$
  (Theorem~\ref{thm:equiv}) contains the seven pairwise disjoint parallel classes
  $\mathcal P_s=\{\beta_{ij}: C(i,j)=s\}$, $s=0,\dots,6$; the remaining $30$ blocks contain
  every point exactly three times.
\item[(iii)] $A$ and $B$ have exactly $14$ common transversals, and at most $7$ of them are
  pairwise disjoint.  So this $4$-GDD has $14$ parallel classes, no eight of them are
  pairwise disjoint, and the $30$ remaining blocks of (ii) cannot be split into three
  parallel classes: the design is three parallel classes short of a $4$-RGDD of type
  $10^4$.
\item[(iv)] $A$ has exactly $1080$ transversals and $305$ orthogonal mates, and for none of
  them, $M$ say, is there a Latin square orthogonal to both $A$ and $M$; so $A$ is not a
  component of a triple, as already known from \cite{MMM07,EW16}; the same holds for $B$
  ($932$ transversals, $5$ mates) and $C$ ($864$ transversals, $4$ mates).  Over all pairs
  of mates of $A$, $91$ is the largest number of distinct ordered pairs.
\end{enumerate}
\end{proposition}

\begin{proof}
(i)--(iii) are checked by \texttt{verify\_10\_4.py} from the definitions (for (iii) it
enumerates all common transversals and finds a largest set of pairwise disjoint ones).
For (ii): $C\perp A$, and the classes of the symbols $0,\dots,6$ of $C$ contain no repeated
pair, so they are common transversals of $A$ and $B$; each point lies in $10$ blocks, seven
of them in these classes.  (iv) is the complete test of Section~\ref{sec:search}.
\end{proof}

In the language of designs, this is the best known approximation to a $4$-RGDD of type
$10^4$ that we are aware of: a $4$-GDD of type $10^4$ with seven pairwise disjoint
parallel classes (Brown, Hedayat and Parker \cite{BHP93} had four).

\subsection{A pair of orthogonal Latin squares from our search}\label{sec:ourpair}

The best pair found by our search (Section~\ref{sec:search}) is shown in
Figure~\ref{fig:ourpair}: two orthogonal Latin squares $A,B$ of order $10$ and, in the
third array, 3 pairwise disjoint common transversals of $A$ and $B$ (cells marked
$0,\dots,2$).

\begin{figure}[ht]
\centering
\setlength{\arraycolsep}{2.6pt}
$\begin{array}{*{10}{c}@{\hspace{2em}}*{10}{c}@{\hspace{2em}}*{10}{c}}
7 & 2 & 1 & 4 & 5 & 8 & 6 & 9 & 0 & 3 & 5 & 6 & 2 & 4 & 1 & 7 & 3 & 8 & 9 & 0 & \cdot & \cdot & \cdot & \cdot & 0 & 1 & \cdot & \cdot & 2 & \cdot\\
3 & 4 & 5 & 0 & 8 & 6 & 9 & 7 & 1 & 2 & 7 & 5 & 9 & 1 & 6 & 2 & 0 & 4 & 8 & 3 & \cdot & \cdot & \cdot & 1 & \cdot & 2 & \cdot & \cdot & \cdot & 0\\
6 & 5 & 8 & 3 & 9 & 2 & 0 & 1 & 7 & 4 & 4 & 7 & 0 & 2 & 9 & 1 & 5 & 3 & 6 & 8 & \cdot & \cdot & 2 & \cdot & \cdot & \cdot & \cdot & \cdot & 0 & 1\\
1 & 3 & 4 & 7 & 2 & 9 & 8 & 0 & 6 & 5 & 9 & 8 & 7 & 3 & 4 & 5 & 2 & 0 & 1 & 6 & 1 & 0 & \cdot & \cdot & \cdot & \cdot & \cdot & \cdot & \cdot & 2\\
9 & 1 & 3 & 8 & 0 & 7 & 5 & 2 & 4 & 6 & 2 & 1 & 6 & 8 & 3 & 9 & 4 & 7 & 0 & 5 & 0 & \cdot & \cdot & \cdot & \cdot & \cdot & 1 & 2 & \cdot & \cdot\\
2 & 0 & 9 & 6 & 7 & 3 & 1 & 4 & 5 & 8 & 0 & 2 & 3 & 7 & 8 & 4 & 6 & 1 & 5 & 9 & \cdot & \cdot & 1 & 0 & 2 & \cdot & \cdot & \cdot & \cdot & \cdot\\
8 & 9 & 2 & 5 & 1 & 0 & 4 & 6 & 3 & 7 & 1 & 4 & 5 & 0 & 7 & 8 & 9 & 6 & 3 & 2 & \cdot & 2 & \cdot & \cdot & \cdot & \cdot & 0 & 1 & \cdot & \cdot\\
5 & 8 & 0 & 1 & 6 & 4 & 7 & 3 & 2 & 9 & 8 & 3 & 4 & 5 & 0 & 6 & 7 & 9 & 2 & 1 & \cdot & \cdot & 0 & 2 & \cdot & \cdot & \cdot & \cdot & 1 & \cdot\\
0 & 7 & 6 & 2 & 4 & 5 & 3 & 8 & 9 & 1 & 6 & 0 & 8 & 9 & 2 & 3 & 1 & 5 & 7 & 4 & \cdot & 1 & \cdot & \cdot & \cdot & \cdot & 2 & 0 & \cdot & \cdot\\
4 & 6 & 7 & 9 & 3 & 1 & 2 & 5 & 8 & 0 & 3 & 9 & 1 & 6 & 5 & 0 & 8 & 2 & 4 & 7 & 2 & \cdot & \cdot & \cdot & 1 & 0 & \cdot & \cdot & \cdot & \cdot\\
\end{array}$
\caption{Orthogonal Latin squares $A$ (left) and $B$ (middle) of order $10$, and 3
pairwise disjoint common transversals of $A$ and $B$ (right).}
\label{fig:ourpair}
\end{figure}

\begin{proposition}\label{prop:ourpair}
For the squares of Figure~\ref{fig:ourpair}, the $100$ blocks $\beta_{ij}$ of
Theorem~\ref{thm:equiv} form a $4$-GDD of type $10^{4}$ with 3 pairwise disjoint parallel
classes (the cells marked $0,\dots,2$).  The squares $A$ and $B$ have exactly
3 common transversals, so this $4$-GDD has exactly 3 parallel classes and is not
resolvable; 7 of the ten pairwise disjoint parallel classes of a resolution are
missing.
\end{proposition}

\begin{proof}
\texttt{verify\_10\_4.py} checks that $A\perp B$, checks the $4$-GDD and the 3 parallel
classes from the definitions, and enumerates all common transversals of $A$ and $B$;
Corollary~\ref{cor:gdd} identifies these with the parallel classes.
\end{proof}



\section{The computer search}\label{sec:search}

\begin{lemma}\label{lem:crit}
A Latin square $L$ of order $n$ is a component of a set of three MOLS if and only if the
cells admit two partitions $\mathcal P$, $\mathcal Q$ into transversals of $L$ with
$|T\cap U|=1$ for all $T\in\mathcal P$, $U\in\mathcal Q$.  Equivalently, for some orthogonal
mate $M$ of $L$ (given by $\mathcal P$) the common transversals of $L$ and $M$, namely the
transversals $U$ of $L$ with $|U\cap T|=1$ for all $T\in\mathcal P$, contain a partition of
the cells.
\end{lemma}

\begin{proof}
If $L,M,M'$ are MOLS take for $\mathcal P,\mathcal Q$ the symbol classes of $M,M'$; the
condition $|T\cap U|=1$ says that $(M,M')$ takes every pair of values once.  Conversely
label the members of $\mathcal P$ and $\mathcal Q$ by $[n]$ to obtain $M$ and $M'$.
\end{proof}

\paragraph{Algorithm.}  For a square $L$ of order $10$ the program
\texttt{mols3\_search.c}
\begin{enumerate}
\item enumerates all transversals of $L$ by backtracking;
\item enumerates \emph{all} orthogonal mates, i.e.\ all exact covers of the $100$ cells
  by transversals (Algorithm~X \cite{Knu00}: branch on an uncovered cell with the fewest
  candidates, prune as soon as some uncovered cell has none);
\item for every mate $\mathcal P$ computes the set $C(\mathcal P)$ of common transversals
  as an intersection of precomputed bit sets $\{U:|U\cap T|=1\}$, $T\in\mathcal P$, and
  decides by a second exact-cover search whether $C(\mathcal P)$ contains a partition of the
  cells;
\item records $\nu(\mathcal P)$, the largest number of pairwise disjoint members of
  $C(\mathcal P)$ (a maximum clique computation), and, for every two mates $M,M'$ of $L$,
  the number of distinct ordered pairs $(M(i,j),M'(i,j))$.
\end{enumerate}
Steps 1--3 decide completely whether $L$ is a component of a triple
(Lemma~\ref{lem:crit}).  By Corollary~\ref{cor:gdd}, $\nu(\mathcal P)$ is the largest number
of pairwise disjoint parallel classes of the $4$-GDD of type $10^{4}$ defined by the pair
$(L,M)$; $\nu=10$ would be a $4$-RGDD of type $10^{4}$.  The pair count in step~4 is the
measure used by Egan and Wanless ($100$ would be a triple, their record is $91$).

\paragraph{Validation.}  The program reproduces the transversal counts of the Cayley
tables of $\Z_7$ ($133$), $\Z_2^{3}$ ($384$), $\Z_3^{2}$ ($2241$), $\Z_9$ ($2025$), and
finds none for $\Z_8$, $\Z_{10}$ and the dihedral group $D_5$ (in accordance with
Proposition~\ref{prop:cm}).  It finds triples of MOLS containing the Cayley tables of
$\Z_7$, $\Z_2^{3}$, $\Z_3^{2}$ and $\Z_9$ (the triples were checked independently).  The
Cayley tables of $\Z_9$ and $\Z_3^{2}$ have $2{,}049{,}219$ and $12{,}445{,}836$ orthogonal
mates respectively (here and below a mate is counted as a partition of the cells into
transversals, i.e.\ up to renaming its symbols), exactly the numbers reported by Egan and
Wanless \cite[Section~5]{EW16}.  An independent Python implementation (dictionary-based
Algorithm~X) reproduced, for $\Z_7$ and for six random squares of order $10$, the
numbers of transversals and mates, the maximum of $|C(\mathcal P)|$ and the number of
mates that extend to a triple ($133$, $635$, $28$, $5$ for $\Z_7$).  Applied to the
squares of Egan and Wanless (Section~\ref{sec:ew}), the C program reproduces the values
$91$ and $7$ reported in \cite{EW16}.

\paragraph{Searches.}
\begin{enumerate}
\item \emph{Random squares.}  Squares sampled with the Jacobson--Matthews Markov chain
  \cite{JM96} (at least $1500$ steps between samples).
\item \emph{Prolongations.}  For every transversal $T$ of the Cayley tables of $\Z_9$ and
  $\Z_3^{2}$, the prolongation to order $10$ (move the symbols on $T$ to a new row and
  column and put the new symbol on $T$); these squares have many more transversals
  ($1376$--$1880$) and up to $2{,}592$ orthogonal mates.
\item A simulated annealing on the number of transversals was also tried; it did not get
  beyond about $980$ transversals (random squares have about $800$) and was abandoned.
\end{enumerate}

\begin{table}[ht]
\centering
\small
\setlength{\tabcolsep}{4pt}
\begin{tabular}{lrrrrrrrr}
\toprule
 & & squares with & & mates with & & & best & \\
family & squares & a mate & mates & $C(\mathcal P)\ne\emptyset$ & $\max|C(\mathcal P)|$ & $\max\nu$ & pairs & triples\\
\midrule
random & 194{,}198 & 117{,}730 & 198{,}229 & 22{,}557 & 3 & 3 & 78 & 0\\
prolongations, $\Z_9$ & 2{,}025 & 1{,}971 & 134{,}136 & 32{,}076 & 1 & 1 & 74 & 0\\
prolongations, $\Z_3^{2}$ & 2{,}241 & 1{,}368 & 264{,}384 & 116{,}640 & 2 & 1 & 75 & 0\\
\bottomrule
\end{tabular}
\caption{Exhaustive per-square tests.  ``mates'': orthogonal mates found (all of them, for every square); ``$C(\mathcal P)\ne\emptyset$'': mates $M$ for which $L$ and $M$ have at least one common transversal; $\nu$: largest number of pairwise disjoint common transversals (disjoint parallel classes of the $4$-GDD of type $10^{4}$ of $(L,M)$; $10$ would be a $4$-RGDD); ``best pairs'': largest number of distinct ordered pairs $(M(i,j),M'(i,j))$ over two mates $M\ne M'$ of one square ($100$ would be a triple).}
\label{tab:results}
\end{table}

\paragraph{Interpretation.}  Typical squares of order $10$ have mates (in our random sample $60.6\%$ have one and the average number of mates is $1.02$, in line with the estimates of \cite{MMM07} quoted in \cite[Section~5]{EW16}: about $60\%$, and an expectation above $1$), but the pairs they
form have almost no common transversals: a $4$-GDD of type $10^{4}$ obtained from such a
pair has at most three disjoint parallel classes in our search, far from the ten needed.  This agrees with
\cite[Table~3]{BKS26}, where none of several hundred orthogonal pairs found by a SAT
solver had more than two common transversals, and is much weaker than the best known
near miss (seven disjoint parallel classes and $91$ distinct pairs for the squares of Egan and
Wanless, Section~\ref{sec:ew}).  These are statistics, not evidence of non-existence: the squares
tested are a vanishing fraction of the $3.48\times 10^{16}$ main classes.

\section{Conclusion}

A $4$-RGDD of type $10^{4}$ is the same object as three MOLS of order $10$
(Theorem~\ref{thm:equiv}).  We could not construct one, and no proof of non-existence is
known; the case $(h,u)=(10,4)$ in \cite{WG14,WG17} is precisely the open problem
``$N(10)\ge 3$?''.  We proved that such a design cannot be obtained from a difference matrix
over a group of order $10$ or by developing base blocks under an abelian point-regular
group (Corollary~\ref{cor:aut}), that a point-regular group would have to be an extension
of $\Z_5$ by the dihedral group of order $8$ (and, by \cite{MMM07}, cannot exist at all),
and that no incomplete design of this kind (such as Brouwer's, with a hole of type $2^{4}$)
can be completed except in the trivial case of a hole of size one
(Proposition~\ref{prop:holes}); in particular a hole of size $4$, which three MOLS of order
$4$ could fill, cannot occur at all (Remark~\ref{rem:h4}).  The closest known approximation, from the squares of Egan and Wanless, is a $4$-GDD of type
$10^{4}$ with seven pairwise disjoint parallel classes; our own search, which tested
$198{,}464$ squares exhaustively, got no further than three disjoint parallel classes (Proposition~\ref{prop:ourpair}).  The
explicit near misses above are verified by \texttt{verify\_10\_4.py}; a candidate triple
found by any other means can be checked with \texttt{python3 verify\_10\_4.py triple FILE},
which builds the $4$-RGDD of Theorem~\ref{thm:equiv} and checks it from the definition.

\begin{thebibliography}{99}

\bibitem{ACD07} R.J.R. Abel, C.J. Colbourn, J.H. Dinitz, Mutually orthogonal Latin squares
(MOLS), in: C.J. Colbourn, J.H. Dinitz (Eds.), Handbook of Combinatorial Designs, 2nd ed.,
Chapman \& Hall/CRC, Boca Raton, 2007.

\bibitem{BCSKG21} C. Bright, K.K.H. Cheung, B. Stevens, I. Kotsireas, V. Ganesh, A SAT-based
resolution of Lam's problem, Proc. AAAI Conf. Artificial Intelligence 35 (2021) 3669--3676.

\bibitem{BKS26} C. Bright, A. Keita, B. Stevens, Myrvold's results on orthogonal triples of
$10\times 10$ Latin squares: a SAT investigation, Electron. J. Combin. 33(1) (2026) \#P1.30;
arXiv:2503.10504.

\bibitem{BS59} R.C. Bose, S.S. Shrikhande, On the falsity of Euler's conjecture about the
non-existence of two orthogonal Latin squares of order $4t+2$, Proc. Nat. Acad. Sci. USA 45
(1959) 734--737.

\bibitem{BSP60} R.C. Bose, S.S. Shrikhande, E.T. Parker, Further results on the
construction of mutually orthogonal Latin squares and the falsity of Euler's conjecture,
Canad. J. Math. 12 (1960) 189--203.

\bibitem{BSH77} A.E. Brouwer, A. Schrijver, H. Hanani, Group divisible designs with block
size four, Discrete Math. 20 (1977) 1--10.

\bibitem{BHP93} J.W. Brown, A.S. Hedayat, E.T. Parker, A pair of orthogonal Latin squares of
order 10 with four shared parallel transversals, J. Combin. Inform. System Sci. 18 (1993)
113--115.

\bibitem{Bro84} A.E. Brouwer, Four MOLS of order 10 with a hole of order 2, J. Statist.
Plann. Inference 10 (1984) 203--205.

\bibitem{Bru63} R.H. Bruck, Finite nets. II. Uniqueness and imbedding, Pacific J. Math. 13
(1963) 421--457.

\bibitem{Del10} E. Delisle, The search for a triple of mutually orthogonal Latin squares of
order ten: looking through pairs of dimension thirty-five and less, M.Sc. thesis,
University of Victoria, 2010.

\bibitem{DH14} P. Dukes, L. Howard, Group divisible designs in MOLS of order ten, Des.
Codes Cryptogr. 71 (2014) 283--291.

\bibitem{EW16} J. Egan, I.M. Wanless, Enumeration of MOLS of small order, Math. Comp. 85
(2016) 799--824.

\bibitem{GW23} M.J. Gill, I.M. Wanless, Pairs of MOLS of order ten satisfying non-trivial
relations, Des. Codes Cryptogr. 91 (2023) 1293--1313.

\bibitem{HP55} M. Hall, L.J. Paige, Complete mappings of finite groups, Pacific J. Math. 5
(1955) 541--549.

\bibitem{JM96} M.T. Jacobson, P. Matthews, Generating uniformly distributed random Latin
squares, J. Combin. Des. 4 (1996) 405--437.

\bibitem{Knu00} D.E. Knuth, Dancing links, in: Millennial Perspectives in Computer Science,
Palgrave, 2000, pp.~187--214.

\bibitem{LTS89} C.W.H. Lam, L. Thiel, S. Swiercz, The non-existence of finite projective
planes of order 10, Canad. J. Math. 41 (1989) 1117--1123.

\bibitem{Man44} H.B. Mann, On orthogonal Latin squares, Bull. Amer. Math. Soc. 50 (1944)
249--257.

\bibitem{MMM07} B.D. McKay, A. Meynert, W. Myrvold, Small Latin squares, quasigroups, and
loops, J. Combin. Des. 15 (2007) 98--119.

\bibitem{Myr99} W. Myrvold, Negative results for orthogonal triples of Latin squares of
order 10, J. Combin. Math. Combin. Comput. 29 (1999) 95--106.

\bibitem{Par59} E.T. Parker, Orthogonal Latin squares, Proc. Nat. Acad. Sci. USA 45 (1959)
859--862.

\bibitem{Par62} E.T. Parker, On orthogonal Latin squares, in: 1960 Institute on Finite
Groups, Proc. Sympos. Pure Math. 6, Amer. Math. Soc., Providence, 1962.

\bibitem{Roy11} D.J. Roy, Confirmation of the non-existence of a projective plane of order
10, M.Sc. thesis, Carleton University, 2011.

\bibitem{Tar1900} G. Tarry, Le probl\`eme des 36 officiers, Compte Rendu de l'Association
Fran\c{c}aise pour l'Avancement des Sciences 29 (1900), part 2, 170--203.

\bibitem{WG14} H. Wei, G. Ge, Some more 5-GDDs, 4-frames and 4-RGDDs, Discrete Math. 336
(2014) 7--21.

\bibitem{WG17} H. Wei, G. Ge, Some more uniformly resolvable designs with block sizes 2 and
4, Discrete Math. 340 (2017) 2243--2249.

\bibitem{ZK15} O. Zaikin, S. Kochemazov, The search for systems of diagonal Latin squares
using the SAT@home project, Int. J. Open Inf. Technol. 3(11) (2015) 4--9.

\bibitem{ZVB25} O. Zaikin, E. Vatutin, C. Bright, Enumerating extended self-orthogonal
diagonal Latin squares of order up to 10, J. Integer Seq. 28 (2025) Article 25.7.4.

\end{thebibliography}

\end{document}
